% ED26-BM-RAT2021-12-1: independently worded complete source treatment.
\chapter{Rationality and Morality: Convergence, Conflict and Guidance}
\label{sec:p11-rationality-12-1}

Fehige and Wessels ask how practical-rational assessments of an action
relate to moral assessments of it. They call this the RM question,
using \(R\) for practical rationality and \(M\) for morality.
Both standards are disputed, as is any claim that one has authority
over the other. Their chapter organizes the debate around sample
desire-based and welfare-based criteria, examines attempts to bring
them together, and defends accepting some divergence rather than
engineering a perfect match at the cost of independently plausible
judgments \parencite[chap.~12.1, pp.~681--692]{FehigeWessels2021Morality}.

The treatment concerns philosophical criteria and arguments, not a
participant study or a numerical estimate of moral conduct. Its scope
is actions and all-things-considered deontic judgments. Beliefs,
desires, emotions, intentions, maxims and whole ways of life may also
be rationally or morally assessed, but the source deliberately does
not undertake those different inquiries here. This action-focused
meaning of practical rationality must not be silently substituted for
every other account in this handbook.

\section{Assessments and the precise convergence question}
\label{sec:p11-rat121-assessment}

\begin{hbdefinition}{def:p11-rat121-deontic}{Required, permitted, optional and forbidden actions}
The source's deontic vocabulary applies separately within each domain.
An action may be rationally required or morally required, rationally
permitted or morally permitted, and so forth. Required actions are
permitted; permitted actions comprise those required and those merely
optional; forbidden actions are outside the permitted class.
\emph{Required}, \emph{obligatory} and \emph{ought to be done}
are the source's roughly synonymous descriptions of the first category.

For this chapter, \emph{rational} means rationally permitted and
\emph{irrational} means rationally forbidden. \emph{Right} means
morally permitted and \emph{wrong} means morally forbidden.
Consequently, calling an act rational or right does not say it is
uniquely obligatory. These conventions reconstruct Figure~12.1.1;
they do not assert an identity between the two domains or supply
a new deontic calculus.
\end{hbdefinition}

The motivating case is an affluent person contemplating spending
80 percent of her income. Morality might require donating it to an
effective lifesaving charity, while practical rationality might require
using it to complete her cherished collection of abstract paintings.
The source asks what follows if both judgments survive examination.
It does not establish either requirement merely from the example's
description, or report how often affluent collectors face it. One must
state the moral criterion, the person's desires and beliefs, the
available alternatives and any relevant incentives before drawing the
two verdicts \parencite[p.~681]{FehigeWessels2021Morality}.

\begin{hbdefinition}{def:p11-rat121-informed}{Well-informedness and the direction of convergence}
An agent is \emph{well-informed}, in the selected source's sense, when
further improvements of her beliefs---correcting false beliefs or
adding true ones---would not change what it is rational for her to do.
The definition concerns stability of the action verdict under better
beliefs, not possession of every truth or infallibility about all topics.

\(R\) and \(M\) \emph{converge} insofar as actions morally
required of well-informed agents are also rationally permitted.
They \emph{diverge} insofar as morally required actions are irrational
for such agents. This is a directed compatibility test. It is not the
claim that every rationally permitted act is morally permitted, that
every morally permitted act is rationally required, or that the
permitted and required classes coincide across domains.
\end{hbdefinition}

Information is isolated because a committed utilitarian can be wrong
about an action's welfare effects. Under a belief-sensitive rational
criterion, following that mistake can be rational yet morally wrong
under an actual-outcome criterion. The well-informedness condition
sets aside that source of difference to investigate more substantive
differences of aim or standard. It neither proves that everyone can
achieve the required information nor supplies a learning algorithm
\parencite[p.~682]{FehigeWessels2021Morality}.

\section{The two sample doctrines and their limitations}
\label{sec:p11-rat121-samples}

\begin{hbdefinition}{def:p11-rat121-instrumentalism}{The source's action-focused instrumentalism}
At the time of action, an available act is rational if and only if
the agent believes no other available act would produce more
fulfillment of her intrinsic desires. An intrinsic desire concerns
something for its own sake, not merely as a means to another end.
The quantity of fulfillment incorporates both how many desires are
fulfilled and their strengths.

The available action, belief and desire are evaluated at the same
time; beliefs and desires can be implicit rather than consciously
entertained. Henceforth, when reconstructing this source, an
unqualified \emph{desire} means an intrinsic desire. The criterion
does not require that such desires be selfish or morally acceptable.
\end{hbdefinition}

The source uses \emph{instrumentalism} for this specific criterion.
It associates beliefs and desire-fulfillment amounts, broadly, with
probabilities and utilities in expected-utility theory. That is an
explanatory comparison, not an exact representation theorem: the
chapter supplies no lottery axioms, probability function, numerical
desire-strength scale or proof that the sample belief criterion is
equivalent to every version of expected-utility maximization.
Note~1 situates instrumentalism and its rivals in Hutcheson,
Millgram, Schroeder, Schmidt and Fehige's earlier discussion.
The mathematical results in P17 retain their own stated premises
\parencite[p.~682 and n.~1]{FehigeWessels2021Morality}.

\begin{hbdefinition}{def:p11-rat121-utilitarianism}{The sample moral and welfare criteria}
The source's sample utilitarian criterion calls an act morally right
if and only if no other available act brings about more welfare
worldwide. Welfare counts regardless of whose it is, and nothing
else is assigned independent moral importance by this sample criterion.

For this comparison the authors additionally suppose that a person's
welfare consists in fulfillment of her desires, with pleasure an
important part through a proposed conceptual connection to desire
fulfillment. This is an adopted welfare theory, not the definition
of every utilitarianism or a proof that every pleasure reduces to
desire fulfillment. Under tied maximizing alternatives more than
one act can be right without each being individually obligatory.
\end{hbdefinition}

The two doctrines disagree in what they optimize and in their
dependence on belief. The first evaluates the act through the
agent's view of her own desires' fulfillment; the second compares
actual welfare worldwide. Well-informedness addresses the belief
difference, but not the possibility that the person's desires favor
paintings while aggregate welfare favors lifesaving. Note~2's
von Wright, Carson and Wessels references concern desire-based welfare;
its Fehige and Heathwood references concern pleasure; Bykvist,
Vaughn and Copp locate rival ethical views. These are external
developments, not imported proofs of the assumptions
\parencite[p.~682 and n.~2]{FehigeWessels2021Morality}.

\section{What makes convergence possible}
\label{sec:p11-rat121-desires}

\subsection{Moral desires, sympathy and a conditional necessity argument}

\begin{hbdefinition}{def:p11-rat121-rm-desires}{RM desires and the direct/causal distinction}
An \emph{RM desire} tends, for a well-informed agent, to make morally
right conduct rational. The qualification is a propensity, not an
assurance that one desire defeats every competing desire.

In the source's Figure~12.1.2, a \emph{moral desire} that \(p\)
directly concerns something morally positive in itself, abbreviated
\(\mathrm{mpi}\). It may invoke morality, such as wanting to do
the right thing or be virtuous, or may instead concern others'
happiness, equal treatment, keeping promises or avoiding lies without
representing those ends as moral. Which of these is \(\mathrm{mpi}\)
depends on the adopted morality.

A \emph{morally causal desire} that \(p\) is indirectly connected
with something \(\mathrm{mpi}\): \(p\) would cause it, would
be caused by it, or would share a cause with it. The kinds can
overlap, even at one classificatory level. This is a classification
of routes to convergence, not a mutually exclusive psychology.
\end{hbdefinition}

Under the sample utilitarianism, Mary can want to do what is right
or want maximal welfare without explicitly thinking of moral status.
Either desire can support convergence when she is well-informed and
its force is not defeated. Sympathy can generate or constitute a
desire that another person fares well. Admiration, indignation or
satisfaction can similarly generate or constitute desires for or
against certain acts. The source locates sympathy-centered ethics
in Schopenhauer and sentiment-centered approaches in Hutcheson and
Hume, with further discussion in Bricke and Fehige and Frank.
These are conceptual and historical comparisons, not measured claims
about universal distributions or causal effects of sentiments
\parencite[pp.~682--683 and nn.~3--4]{FehigeWessels2021Morality}.

The authors ask whether convergent desires are all contingent.
They present an argument from Fehige's 2004 work for a stronger
resource. Its first premise is that fully representing another
person's particular pleasure entails experiencing that pleasure
oneself. Its second premise identifies a disposition to be pleased
on fully representing a state of affairs with a desire that the
state obtain. From those premises, representing another's pleasure
provides a disposition of the specified kind, hence an other-regarding
desire. The proposed strength is proportional to the pleasure represented.
If that conceptual account is sound, the desire need not await an
occurrent act of sympathetic imagining; its dispositional basis
would be present for every other person's pleasure.

The source explicitly presents this as an argument whose success is
conditional. Full representation's phenomenal entailment and the
disposition/desire identification do substantive work; naming the
disposition does not establish the first premise. Representing that
someone enjoys a sensation must not simply be equated with undergoing
it without argument. The proportional-strength and always-present
claims also require that account, not only an instance of empathy.
Even if accepted, these desires can coexist with competing aims;
their existence alone does not show that every morally required
act maximizes the agent's total desire fulfillment. Note~5 explicitly
points to limits in Fehige's separate book, chapter~6. No complete
proof of universal RM convergence is supplied or claimed here
\parencite[p.~683 and n.~5]{FehigeWessels2021Morality}.

\subsection{Why instrumentalism is not psychological egoism}

The source rejects an equivocation: instrumentalism acts through
the agent's own desires; egoism concerns the content of those desires
and the person's actions on them. If the person strongly wants
others to flourish and acts on that desire, acting on her own desire
does not make her an egoist. The argument that instrumentalism is
egoistic can be rescued only by adding that genuinely other-regarding
desires do not exist. The authors report no sound argument for that
addition and cite Hutcheson and Sharp in note~6. This is a conceptual
objection to the inference, not a quantified empirical prevalence claim.

Combining instrumentalism with desire-fulfillment welfare also needs
care. A well-informed Mary at time \(t\) can maximize her current
desires' fulfillment and, on the adopted welfare theory, her current
person-stage's welfare. She can nevertheless be motivated entirely
by concern for others. The current-stage result does not entail
maximal welfare for her entire temporally extended person, whose
later desires can differ. If welfare is instead identified with
experienced pleasure while instrumentalism is retained, current
desire fulfillment no longer establishes even that limited welfare
identity \parencite[pp.~683--684]{FehigeWessels2021Morality}.

The RM question consequently differs from reconciling egoism,
prudence, self-interest, self-love or personal welfare with morality.
The source mentions Plato's invisible Gyges, Aristotle's virtue and
flourishing, and Sidgwick's individual/universal happiness as important
historical discussions that answer the broader RM question only
partly or indirectly. Note~7 connects Kavka's wider reconciliation
project and Cholbi's rational-egoist discussion. Those historical
cases do not justify making selfish content part of instrumentalism's
definition.

\subsection{Causally connected desires and incentives}

All three indirect routes in Figure~12.1.2 matter. Wanting to
cultivate a garden, clean a kitchen or mend bones can produce pleasure
for beholders, diners or patients: the desired activity causes the
morally positive result. Mary can want cake that happy parents give
her because they are happy: something morally positive causes the
desired result. Mary can want praise at a charity ball while being
indifferent to malaria victims: her donation causes both praise
and saved lives, without either effect causing the other.
The source's third example is specifically a common-cause route,
not a claim that praise itself saves the victims
\parencite[p.~684]{FehigeWessels2021Morality}.

External honor, help, exclusion or punishment can link morally
positive behavior to desired consequences. Internal joy, remorse
or their avoidance can do so too. Wanting to do good and wanting
the satisfaction of having done good are distinct desires, as
note~8's Hutcheson, Rashdall and Schlick discussion emphasizes.
The same desire for one's joy can belong to both kinds: its joy
may be caused by a good act, and may itself be morally positive
on the adopted welfare view. The causal and intrinsic-value grounds
for the two classifications remain distinct.

The source surveys game-theoretic and agent-based accounts of
cooperation under incentives, especially repeated encounters with
rewards and punishments. Note~9 identifies Fudenberg and Tirole and
Maschler, Solan and Zamir for folk-theorem approaches, and Gotts,
Polhill and Law and Hegselmann for simulations. It supplies no
single theorem's hypotheses or proof and no simulated population,
parameter set, outcome measure or uncertainty estimate. These are
precisely identified survey boundaries, not a locally established
theorem that repetition always creates cooperation, justice or
morally optimal outcomes. A game payoff is not automatically a
moral value \parencite[pp.~684--685 and n.~9]{FehigeWessels2021Morality}.

Convergent incentives and moral desires can therefore explain why
right conduct is often rational without establishing full convergence.
In the source's collector example, if the person's competing desires
and incentives favor paintings sufficiently, the two sample criteria
still conflict. Neither knowing the actual effects nor recognizing
the moral verdict alone supplies the missing motivational ranking.

\section{Changing rationality to resist divergence}
\label{sec:p11-rat121-moralizing}

\begin{hbdefinition}{def:p11-rat121-rationalism}{Moral rationalism and two directions of revision}
The source calls the claim that every morally required act is rational
\emph{moral rationalism}; note~10 also records the overridingness
terminology. One can try to secure convergence by changing practical
rationality toward the moral standard, by reducing or reshaping moral
requirements toward practical rationality, or by changing both.
Assuming moral rationalism in order to select the two standards is
a substantive commitment, not an independent proof of their fit.
\end{hbdefinition}

Note~10 points to Scheffler, Cholbi, Portmore, Dorsey and Brink for
the dispute. Fehige and Wessels next distinguish relatively modest
departures from their instrumentalist sample from more fundamental
changes in its grounds \parencite[p.~685]{FehigeWessels2021Morality}.

\subsection{Temporal scope, collective traps and dispositions}

Broadening the conative basis to include past, present and future
desires can reduce an agent's indifference toward her earlier or
later selves. This is an intrapersonal improvement, not an
other-regarding bridge. Concern for all one's person-stages does
not entail concern for everyone's welfare. The source does not
provide a complete aggregation rule for conflicting desires across
time, so temporal broadening is not a specified maximization algorithm.

A different motivation comes from traps of practical rationality:
if all agents act rationally under their individual criteria, each
can obtain less desire fulfillment than under a collectively different
course of action. The source's arms-race illustration has each person
expect to benefit from carrying a weapon regardless of whether others
do. The collective outcome is a worse armed society than the
weaponless alternative. Sanctions could alter individual incentives,
but administering them also uses resources. This is a stipulated
strategic illustration, not an empirical estimate that every society's
weapon choices have these payoffs. Note~11 points to Peterson,
Dawes and Binmore for the broader social-dilemma literature.

The source surveys Gauthier's constrained maximization as a response.
An agent can evaluate adopting a recognizably cooperative disposition
by its utility benefits, then let that disposition constrain later
choices rather than maximize independently at each moment. The
disposition contributes to a mutually better outcome conditional on
others' corresponding contribution. Recognition can open cooperative
opportunities unavailable to an unconstrained maximizer. If those
claims hold, choosing the disposition can be rational even though
some later acts would not be chosen by isolated act maximization.
This distinguishes evaluating a disposition from changing the payoff
ranking of an unchanged one-shot act.

Fehige and Wessels regard the proposal as controversial and bounded
by mutual advantage: it does not aim to establish duties toward
beings who need help but offer no return. Recognition, partner
responses and the available alternatives are essential dependencies,
not guaranteed by calling someone constrained. Note~12 compares
McClennen's resolute choice and records that Gauthier's 2013 conception
is acknowledged to be morally charged. It also names symmetry/mirror
arguments, dependency equilibria, correlated-equilibrium antecedents,
metagames with partial commitments and team reasoning as distinct,
more complex approaches. Their full procedures and proofs are not
supplied by this survey \parencite[pp.~685--686 and nn.~11--12]{FehigeWessels2021Morality}.

\subsection{Worldly reasons rather than desires}

A larger departure makes action rational when the agent believes
there is no stronger reason to perform another available action.
The source illustrates facts about pains and cures and proposes
desire-independent reasons: completing Peter's education can count
for donating even when the donor wants neither that completion nor
anything connected with it. The deliberately different moon-serenading
example tests what a proposed realm of worldly reasons can contain.
These are candidates within a normative theory, not established
medical advice or observations of desires.

\begin{hbwarning}{def:p11-rat121-worldly-bridge}{Three independent dependencies of a reasons route}
To obtain convergence through worldly reasons, the reasons must
exist; their comparative force must have the relevant affinity
with the moral requirements; and the proposed connection from
beliefs about reasons to rational action must be defensible.
Replacing \emph{desire} with \emph{reason} establishes none
of those three claims.

If the strongest relevant reasons always favor maximizing global
welfare, the route can support the sample utilitarian criterion.
If reasons instead support promises, loyalty or knowledge in the
proportions required by another morality, it can support that
morality. In each case it is the substantive comparative-reasons
premise that does the work, not the mere existence of some reason.
\end{hbwarning}

Notes~13--14 identify distinct belief-about-reasons formulations in
Scanlon and Parfit, and worldly-reasons approaches in Nagel,
Scanlon, Dancy, Foot and Stoutland. Their differences should not
be erased by the shared language. The criterion here concerns
what the agent believes about comparative reasons, not immediate
identity of rationality with successful response to all actual reasons
\parencite[p.~686 and nn.~13--14]{FehigeWessels2021Morality}.

\subsection{Autonomy and the Kantian route}

The source also interprets Kant as defending concern for other
rational beings' ends, including advancing them as far as possible.
This is the source's rational-benevolence interpretation, not an
identification of Kantian ethics with the sample welfare utilitarianism.
Its proposed connection runs through autonomy:
the rational will gives itself law rather than being governed by
inclinations, and the resulting respect concerns rationality,
rational beings and their ends generally rather than one privileged
person or group. That is meant to connect rational agency with
benevolence without appealing to the agent's antecedent desires.

The authors stress how contentious the connection is, even among
Kant's followers, and do not claim to have established its missing
justificatory steps. Freedom from inclination does not by itself
prove the particular positive duties or the proposed generality.
Note~15 locates the \emph{Groundwork} passages and their translation;
it also records the contested role of the later \emph{Critique}'s
fact of reason---whether replacement, summary or supplement of
the earlier justification---and Smith's later rationality-respect
approach. Those historically different routes remain distinct
\parencite[p.~686 and n.~15]{FehigeWessels2021Morality}.

\section{Changing morality to resist divergence}
\label{sec:p11-rat121-rationalizing}

The converse strategy reduces the demands that conflict with agents'
aims. The authors contrast utilitarian impartiality with the partiality
of many people's projects and attachments. This is their philosophical
motivation, not a measured statement about the proportion of all
agents whose desire profiles satisfy impartiality.

\begin{hbdefinition}{def:p11-rat121-prerogatives}{Three source-specific agent prerogatives}
The source describes permission to give one's welfare up to
\(k\) times the weight assigned to others' welfare, for \(k>1\);
permission never to surrender more than \(l\) units of one's
welfare, for \(l>0\); and permission never to let one's welfare
fall below a threshold \(m\). These are three distinct forms,
not one formula. None supplies a common interpersonal welfare
scale, a universal parameter value or a moral justification by itself.

Other prerogatives can concern projects or people close to the agent.
They permit partiality even where the worldwide welfare total
would be lower. Note~16 associates the three forms respectively
with Scheffler, Mulgan and Nagel, and names Wessels and Stroud
for further variants.
\end{hbdefinition}

Relaxing requirements can narrow the gap without closing it. A
remaining demand to consider another's interests can still conflict
with some fully informed desire profile. One must separately
establish why those remaining demands are always rational rather
than infer convergence from their moderation. The source notes that
some prerogativists adopt partiality precisely to maintain the desired
relation to rationality; that is an acknowledged motivation for
revision, not an independent convergence argument
\parencite[pp.~686--687 and n.~16]{FehigeWessels2021Morality}.

Hobbesian reciprocitarian contract theories occupy a related middle
position in the source's comparison. Their rules make each party's
desire fulfillment greater under mutual compliance than under
noncompliance by all. Mutual restraint from insulting each other is
the simple example. Limiting obligations to mutual benefit reduces
their extent, particularly for beings unable to offer a reciprocal
advantage. This is the specified reciprocitarian rationale, not
every possible theory of contract or every institution protecting
dependents.

Joint advantage does not establish that each isolated act of
compliance maximizes the actor's desire fulfillment. The authors
allow situations in which a person can profit from breaking a
rule and escaping sanctions. They take actual sanctioning regimes
to leave gaps; the chapter does not provide a formal theorem or
a measured rate of undetectable wrongdoing. Note~17 identifies
Trapp, Kavka and the later Stemmer as acknowledging rational/moral
divergence, and distinguishes avoiding it through Gauthier's changed
rationality criterion. Changing the contract's content and changing
the criterion of rational action are different revisions
\parencite[p.~687 and n.~17]{FehigeWessels2021Morality}.

The radical option combines moral rationalism with a broadly
instrumentalist basis for reasons. In the Harman account the
authors discuss, asserting that a person morally ought to act
implies that the person has a relevant reason, and such reasons
require appropriate desires, goals or intentions. If those attitudes
are absent, the proposed moral ought is withdrawn rather than
allowed to conflict with practical rationality. Harman's examples
include an alien without concern for us and Hitler's atrocities.
Fehige and Wessels emphasize the moral cost: this view denies
the relevant moral-wrongness or moral-ought judgment where ordinary
condemnation would insist on it, though other condemnatory terms
remain available within Harman's view. Their reconstruction reports
the position; it does not endorse excusing atrocities.

Convergence is obtained here by restricting the addressees or
instances of the moral ought. The authors call that strategy
Procrustean: the fit follows because the misfitting demands were
excluded. The fitting procedure still owes a defense of that moral
restriction. Note~18 attributes the claims to Harman 1975/1977
and records that Stemmer retracted a similar stronger sanction-bound
position in 2017. A source-specific historical account must not
attribute the retracted 2008 view as his unchanged position
\parencite[p.~687 and n.~18]{FehigeWessels2021Morality}.

\section{Accepting divergence without giving up guidance}
\label{sec:p11-rat121-divergence}

\subsection{Independent plausibility before guaranteed harmony}

The authors recommend accepting divergence when the independently
plausible accounts suggest it and no adequate reason requires
theoretical resistance. The desire-based view has intuitive force
because a person's actual aims seem relevant to what it is rational
for her to do. The moral view has independent force because
indifference to others and freedom from backlash do not ordinarily
seem to permit harming them. Identifying the standards wholesale
can sacrifice one of these judgments. Meeting halfway merely
distributes that cost unless the revised features have independent
support. The claim is an argued philosophical assessment, not a
proof that no conceivable compatible theory could be developed
\parencite[pp.~687--688]{FehigeWessels2021Morality}.

\subsection{Four objections and the source's replies}

The first objection appeals to bad noncompliance: if people generally
act rationally, making a moral requirement irrational makes it
unlikely to be fulfilled, with harmful consequences. The reply
separates a sad consequence from evidence that a theory is false.
Badness may justify efforts to alter incentives or desires, but
does not establish that the rational or moral classification was
mistaken. The general compliance premise is not an empirical
rate established by this chapter.

The second objection calls moral judgment pointless when it cannot
motivate the addressee to comply. The source replies that moral
thinking and speech have many functions and indirect effects.
Its example considers judging that Mary ought to spend half her
spare time fighting for a cause. Even if Mary never does, the
judgment can help form decisions, education, relationships,
politics or sanctions, including those of other people. This is
an account of possible functions, not a measured intervention
effect or a claim that every moral utterance is beneficial.

The third objection identifies justification of morality with
showing that compliance would be rational. The authors answer
that justification is addressed to someone. Even on that account,
failure to justify the requirement to every person it addresses
does not entail failure to justify it to anyone. Moreover, the
proposed conception of moral justification is itself controversial,
and justification must eventually reach starting points not further
justified in that particular argument. These are distinct replies:
the limited quantifier, the disputed criterion, and the status
of justificatory starting points. The last does not license using
any unsupported premise merely because an argument stops somewhere.

The fourth objection invokes normative impasses: a well-informed
agent faces an irrational moral requirement and has no third
standard to adjudicate. A theory leaving her with incompatible
directions seems unhelpful. The source replies first that discomfort
does not show the theory misdescribes the norms; independent
evidence of a mistake is needed. Second, divergence does not
entail absence of adjudication. One domain might prevail, or
an overarching norm might settle the conflict while the domain-
specific judgments still differ. Note~19 distinguishes pluralist
accounts with overall authority from accounts denying an overall
ought. It names McLeod, Woods and Case for adjudication approaches,
and Baker for skepticism. None is made mandatory merely by the
existence of conflict \parencite[pp.~688--690 and n.~19]{FehigeWessels2021Morality}.

\subsection{The all-practical-reasons reply and its remaining assumptions}

The source then develops its own guidance reply. Practical reason
in a comprehensive sense considers all reasons to do things.
First distinguish a failure of a verdict within practical reason
from a conflict between it and morality. If reasons themselves
are incomparable or yield no overall balance, lack of guidance
already exists within the practical domain. Excluding RM divergence
would not by itself remedy that distinct problem.

Consider instead the case in which the balance of all practical
reasons gives something the agent has most reason to do. Moral
practical reasons can then be interpreted in two ways. If they
really are practical reasons of a particular kind, the inclusive
balance already takes them into account, as yellow bicycles are
included among bicycles. If the phrase does not denote practical
reasons at all, their exclusion from that balance is no omission
of a reason to act, as toy money need not be money. In either
case the stipulated overall balance supplies practical guidance.
This does not show that every moral requirement wins that balance;
it shows why a difference of domain judgments need not destroy it
\parencite[pp.~688--689]{FehigeWessels2021Morality}.

The conditions are important. The reply assumes a comprehensive
classification of practical reasons and, in the selected case,
an available comparative verdict. It does not prove that every
situation has a unique strongest action, give an aggregation
algorithm, or justify excluding moral considerations by fiat.
Nor is its comprehensive use of \emph{practical reason} an
unannounced derivation that the earlier instrumentalist sample
already includes every worldly normative reason. Different
conceptions of that domain can agree about the structural reply
while disagreeing over its reasons and verdict.

The chapter's conclusion is therefore neither that morality is
pointless nor that rationality and morality are identical. It
distinguishes two shortcomings: an act can be instrumentally
irrational yet benevolent and right, or instrumentally rational
yet malevolent and wrong. The remaining theoretical inquiry
concerns the kinds, mechanisms and extent of their relation.
The practical recommendation to shape desires and their means
of fulfillment so as to reduce conflict is expressly a moral
ought in the source, not a theorem following from observed
cooperation \parencite[p.~689]{FehigeWessels2021Morality}.

\section{Neighboring accounts and a consultation procedure}
\label{sec:p11-rat121-consultation}

Broome's Chapter~\ref{sec:p11-reasons-and-rationality} asks whether
rationality can be reduced to correct response to actual or believed
reasons. His mental/structural conception is not Fehige and Wessels'
stipulated desire-fulfillment criterion. His distinction between
a correctness standard and its normative authority also prevents
the present moral/rational classifications from settling an
overarching ought merely by their names. The worldly-reasons route
here introduces three controversial dependencies; it does not
refute Broome's objections by relabeling successful compliance rational.

TSR in Chapter~\ref{sec:p11-rationality-10-5} considers intended
patterns, continuing commitments, multiple practical reasons and
interpersonal practices. Its nonconsequentialist promise or role
reasons are a substantive alternative to pointwise desire fulfillment.
This relates to the present temporal and collective revisions but
is not identical to either. TSR's acceptance within a form of life,
Gauthier's advantage from a cooperative disposition and the present
claim about the strongest practical reason supply different
arguments and limitations. An accepted practice, mutual advantage
and moral rightness cannot be equated without those arguments.

Adaptationism in Chapter~\ref{sec:p11-rationality-10-6} evaluates
evolved mechanisms relative to selected ancestral functions. Its
explanation of cooperation under uncertain repetition is not a
proof that every act favored by fitness is morally required, nor
that all rational action is based on the agent's intrinsic desires
at the moment of acting. The present morality question remains
even when a behavior is prevalent, stable or well designed for
an adaptive problem. Evolutionary payoffs and moral welfare
also need separate interpretations in P17's population models.

Wedgwood's Chapter~\ref{sec:p11-rationality-2-2} distinguishes
attitudes and rational formation from external correctness.
Here the action/outcome moral criterion and belief-sensitive
rational criterion supply one specific contrast, not a unification
of theoretical and practical rationality. P19's
Definition~\ref{def:p19-benchmark} separates normative comparison
from observed response; neither a sympathy report nor conformity
to a sanction establishes a moral requirement.

\begin{hbmethod}{def:p11-rat121-use}{Locating the actual RM claim before comparing verdicts}
Identify the action, time, alternatives and the agent's beliefs
and intrinsic desires. State the rational criterion and the
moral criterion independently, including the welfare interpretation
or other moral grounds. Distinguish required, permitted, optional
and forbidden verdicts; do not replace the source's directed
convergence claim by identity of all verdicts.

Ask whether improved beliefs would change the rational verdict.
Then identify any direct moral desire or causal incentive, its
strength relative to competing desires, and the actual causal link.
Test the charity-praise case: common cause is not the same as
desiring victims' welfare or praise causing their rescue.

If convergence is secured by revision, identify which standard
changed and why the change is independently plausible. Temporal
extension does not create interpersonal concern; a cooperative
disposition is not act maximization in an unchanged dilemma;
modest prerogatives are not elimination of all moral demands.
If invoking worldly reasons, supply existence, comparative moral
fit and the rationality bridge separately.

If divergence remains, distinguish it from lack of guidance.
State the proposed overriding or adjudicative standard and its
authority rather than inventing one to make the result comfortable.
Specify whether the all-reasons balance actually gives a verdict;
incomparability within that balance is another problem.
Empirical claims need the separate task, population, design,
measure, uncertainty and transfer record of
Method~\ref{def:p19-evidence-record}.
\end{hbmethod}

A successful application can identify the same well-informed
charitable act as right and rational because direct concern for
others or the stated incentives actually favors it. A failure
case is the knowledgeable collector whose sufficiently stronger
desire for paintings conflicts with the lifesaving requirement.
Another failure is claiming that the conflict vanished while
quietly dropping the obligation to a powerless beneficiary.
These examples make different premises visible rather than
pretend to settle every disputed view of rationality or morality.

\index{rationality!morality and convergence}
\index{moral rationalism}
\index{instrumentalism!practical rationality}
\index{moral desires}
\index{morally causal desires}
\index{normative guidance!divergent domains}
\index{agent prerogatives}
